\documentclass[11pt,a4paper]{article}
\usepackage[margin=2.5cm]{geometry}
\usepackage{amsmath,amssymb,amsthm}
\usepackage{physics}
\usepackage{booktabs}
\usepackage{hyperref}
\usepackage{array}
\usepackage{tcolorbox}
\usepackage{tikz}
\usetikzlibrary{arrows.meta,patterns}

\newcommand{\Rxi}{R_\xi}
\newcommand{\Mpl}{M_{\mathrm{Pl}}}
\newcommand{\Mplbar}{\bar{M}_{\mathrm{Pl}}}
\newcommand{\Mfive}{M_5}
\newcommand{\calI}{\mathcal{I}}
\newcommand{\GN}{G_N}
\newcommand{\MZ}{M_Z}
\newcommand{\mgap}{m_{\mathrm{gap}}}

\title{Derivation v22: KK Conventions Unification\\
Interval vs Circle vs Orbifold and Canonical $\Rxi$}
\author{EDC Collaboration\\[3pt]
\small\textit{Derivation note; convention harmonization for BLOCK-003}}
\date{February 2026}

\begin{document}
\maketitle

\begin{abstract}
This note resolves the $\pi$-factor discrepancy between derivations v15--v20
(which used $\Rxi = \hbar c/\MZ$) and v21 (which derived $\Rxi = \pi\hbar c/\MZ$
from the KK mass gap). We derive the KK spectral quantization for three standard
geometries: (A) flat interval with Neumann BCs, (B) circle $S^1$, and
(C) orbifold $S^1/\mathbb{Z}_2$. A conventions dictionary maps between
parametrizations and shows that the $\pi$ difference is purely definitional:
it depends on whether $\Rxi$ denotes the interval length $L$ or the circle
radius $R = L/\pi$. We establish a canonical convention for BLOCK-003 going forward.
\end{abstract}

%═══════════════════════════════════════════════════════════════════════════════
\section{The Problem: Convention Clash}
%═══════════════════════════════════════════════════════════════════════════════

In the BLOCK-003 gravity program, two conventions for $\Rxi$ have emerged:

\textbf{Convention I (v15--v20):}
\begin{equation}
\Rxi^{(\mathrm{old})} = \frac{\hbar c}{\MZ} = 2.165 \times 10^{-18}~\mathrm{m}
\label{eq:Rxi-old}
\end{equation}

\textbf{Convention II (v21):}
\begin{equation}
\Rxi^{(\mathrm{new})} = \frac{\pi \hbar c}{\MZ} = 6.80 \times 10^{-18}~\mathrm{m}
\label{eq:Rxi-new}
\end{equation}

The ratio is:
\begin{equation}
\frac{\Rxi^{(\mathrm{new})}}{\Rxi^{(\mathrm{old})}} = \pi
\label{eq:ratio-pi}
\end{equation}

This propagates to $\Mfive$ via $\Mfive = \Mplbar^{2/3}\Rxi^{-1/3}$:
\begin{equation}
\frac{\Mfive^{(\mathrm{old})}}{\Mfive^{(\mathrm{new})}} = \pi^{1/3} \approx 1.46
\label{eq:M5-ratio}
\end{equation}

\textbf{Question:} Is this a physical difference or a definitional ambiguity?

\textbf{Answer (this note):} It is \emph{purely definitional}. The $\pi$ arises from
whether $\Rxi$ is defined as the interval length $L$ or the circle radius $R = L/\pi$.
Both conventions describe the same physics; the choice is conventional.

%═══════════════════════════════════════════════════════════════════════════════
\section{Definitions and Notation}
%═══════════════════════════════════════════════════════════════════════════════

\begin{tcolorbox}[title=Master Definitions, colback=blue!5, colframe=blue!50]
\textbf{Geometric parameters:}
\begin{itemize}
\item $L$ = length of the interval $[0, L]$
\item $R$ = radius of a circle $S^1$ with circumference $2\pi R$
\item $\Rxi$ = the compactification scale used in BLOCK-003 (to be fixed)
\end{itemize}

\textbf{Coordinate domains:}
\begin{itemize}
\item Interval: $\xi \in [0, L]$
\item Circle: $\xi \in [0, 2\pi R)$ with periodic identification $\xi \sim \xi + 2\pi R$
\item Orbifold $S^1/\mathbb{Z}_2$: fundamental domain $\xi \in [0, \pi R]$ with branes at fixed points
\end{itemize}

\textbf{Boundary conditions (BCs):}
\begin{itemize}
\item Neumann (N): $\partial_\xi \psi = 0$ at boundary
\item Dirichlet (D): $\psi = 0$ at boundary
\item Periodic (P): $\psi(\xi + 2\pi R) = \psi(\xi)$
\end{itemize}

\textbf{Mass gap:} $\mgap \equiv m_1$ = mass of first KK excitation above zero mode
\end{tcolorbox}

%═══════════════════════════════════════════════════════════════════════════════
\section{Case A: Flat Interval with Neumann--Neumann BCs}
%═══════════════════════════════════════════════════════════════════════════════

\subsection{Setup}

Consider a flat 5D spacetime with metric:
\begin{equation}
ds^2 = \eta_{\mu\nu} dx^\mu dx^\nu + d\xi^2, \quad \xi \in [0, L]
\label{eq:caseA-metric}
\end{equation}

The 5D graviton mode equation (from linearized Einstein equations) for a flat
background is:
\begin{equation}
\psi_n''(\xi) + m_n^2 \psi_n(\xi) = 0
\label{eq:caseA-mode-eq}
\end{equation}

General solution:
\begin{equation}
\psi_n(\xi) = A_n \cos(m_n \xi) + B_n \sin(m_n \xi)
\label{eq:caseA-general}
\end{equation}

\subsection{Applying Neumann BCs}

At $\xi = 0$:
\begin{equation}
\psi_n'(0) = B_n m_n = 0 \quad \Rightarrow \quad B_n = 0 \text{ (for } m_n \neq 0\text{)}
\label{eq:caseA-BC-left}
\end{equation}

Solution reduces to:
\begin{equation}
\psi_n(\xi) = A_n \cos(m_n \xi)
\label{eq:caseA-cos-form}
\end{equation}

At $\xi = L$:
\begin{equation}
\psi_n'(L) = -A_n m_n \sin(m_n L) = 0
\label{eq:caseA-BC-right}
\end{equation}

For non-trivial $A_n \neq 0$:
\begin{equation}
\sin(m_n L) = 0 \quad \Rightarrow \quad m_n L = n\pi, \quad n = 0, 1, 2, \ldots
\label{eq:caseA-quantization}
\end{equation}

\subsection{Spectrum (Case A)}

\begin{equation}
\boxed{m_n = \frac{n\pi}{L}, \quad n = 0, 1, 2, \ldots}
\label{eq:spectrum-A}
\end{equation}

\textbf{Zero mode:} $m_0 = 0$ with $\psi_0 = \mathrm{const}$ (massless 4D graviton).

\textbf{Mass gap:}
\begin{equation}
\mgap^{(A)} = m_1 = \frac{\pi}{L}
\label{eq:mgap-A}
\end{equation}

\textbf{Inversion:}
\begin{equation}
\boxed{L = \frac{\pi}{\mgap}}
\label{eq:L-from-mgap}
\end{equation}

\subsection{Normalized Mode Functions}

With normalization $\int_0^L |\psi_n|^2 d\xi = 1$:
\begin{align}
\psi_0 &= \frac{1}{\sqrt{L}} \label{eq:caseA-psi0}\\
\psi_n &= \sqrt{\frac{2}{L}} \cos\left(\frac{n\pi\xi}{L}\right), \quad n \geq 1
\label{eq:caseA-psin}
\end{align}

%═══════════════════════════════════════════════════════════════════════════════
\section{Case B: Circle $S^1$ with Periodic BCs}
%═══════════════════════════════════════════════════════════════════════════════

\subsection{Setup}

Consider a circle with radius $R$ and circumference $C = 2\pi R$:
\begin{equation}
ds^2 = \eta_{\mu\nu} dx^\mu dx^\nu + d\xi^2, \quad \xi \in [0, 2\pi R)
\label{eq:caseB-metric}
\end{equation}

with periodic identification:
\begin{equation}
\xi \sim \xi + 2\pi R
\label{eq:caseB-periodic}
\end{equation}

\subsection{Mode Equation and Periodicity}

The mode equation is the same:
\begin{equation}
\psi_n''(\xi) + m_n^2 \psi_n(\xi) = 0
\label{eq:caseB-mode-eq}
\end{equation}

General solution:
\begin{equation}
\psi_n(\xi) = A_n e^{im_n\xi} + B_n e^{-im_n\xi}
\label{eq:caseB-general}
\end{equation}

Periodicity condition $\psi_n(\xi + 2\pi R) = \psi_n(\xi)$ requires:
\begin{equation}
e^{im_n \cdot 2\pi R} = 1 \quad \Rightarrow \quad m_n \cdot 2\pi R = 2\pi n
\label{eq:caseB-quantization-step}
\end{equation}

\subsection{Spectrum (Case B)}

\begin{equation}
\boxed{m_n = \frac{n}{R}, \quad n = 0, \pm 1, \pm 2, \ldots}
\label{eq:spectrum-B}
\end{equation}

\textbf{Note:} Negative $n$ gives the same mass $|m_n|$; there are two independent
modes for each $|n| \geq 1$ (complex conjugates or $\sin/\cos$).

\textbf{Zero mode:} $m_0 = 0$ (massless graviton).

\textbf{Mass gap:}
\begin{equation}
\mgap^{(B)} = m_1 = \frac{1}{R}
\label{eq:mgap-B}
\end{equation}

\textbf{Inversion:}
\begin{equation}
\boxed{R = \frac{1}{\mgap}}
\label{eq:R-from-mgap}
\end{equation}

\subsection{Alternative: Circumference Parametrization}

Define $C = 2\pi R$ (circumference). Then:
\begin{equation}
m_n = \frac{2\pi n}{C}
\label{eq:spectrum-B-circ}
\end{equation}

and the mass gap is:
\begin{equation}
\mgap^{(B)} = \frac{2\pi}{C} \quad \Rightarrow \quad C = \frac{2\pi}{\mgap}
\label{eq:C-from-mgap}
\end{equation}

\subsection{Normalized Mode Functions}

With normalization $\int_0^{2\pi R} |\psi_n|^2 d\xi = 1$:
\begin{align}
\psi_0 &= \frac{1}{\sqrt{2\pi R}} \label{eq:caseB-psi0}\\
\psi_n &= \frac{1}{\sqrt{\pi R}} \cos\left(\frac{n\xi}{R}\right), \quad n \geq 1
\label{eq:caseB-psin-cos}\\
\tilde{\psi}_n &= \frac{1}{\sqrt{\pi R}} \sin\left(\frac{n\xi}{R}\right), \quad n \geq 1
\label{eq:caseB-psin-sin}
\end{align}

%═══════════════════════════════════════════════════════════════════════════════
\section{Case C: Orbifold $S^1/\mathbb{Z}_2$}
%═══════════════════════════════════════════════════════════════════════════════

\subsection{Setup}

Start with circle $S^1$ of radius $R$ and impose $\mathbb{Z}_2$ identification:
\begin{equation}
\xi \sim -\xi
\label{eq:caseC-Z2}
\end{equation}

The fundamental domain is:
\begin{equation}
\xi \in [0, \pi R]
\label{eq:caseC-domain}
\end{equation}

Fixed points (branes) are at $\xi = 0$ and $\xi = \pi R$.

\subsection{Even Modes (Graviton)}

Graviton modes must be even under $\mathbb{Z}_2$:
\begin{equation}
\psi_n(-\xi) = \psi_n(\xi)
\label{eq:caseC-even}
\end{equation}

This requires the $\cos$ solution:
\begin{equation}
\psi_n(\xi) = A_n \cos(m_n \xi)
\label{eq:caseC-cos}
\end{equation}

Combined with periodicity on the covering space ($\psi_n(\xi + 2\pi R) = \psi_n(\xi)$):
\begin{equation}
\cos(m_n(\xi + 2\pi R)) = \cos(m_n \xi) \quad \Rightarrow \quad m_n \cdot 2\pi R = 2\pi n
\label{eq:caseC-periodicity}
\end{equation}

\subsection{Spectrum (Case C, Even Modes)}

\begin{equation}
\boxed{m_n = \frac{n}{R}, \quad n = 0, 1, 2, \ldots}
\label{eq:spectrum-C}
\end{equation}

\textbf{Same as circle!} But now only even $n$ modes are present (no double degeneracy).

\textbf{Mass gap:}
\begin{equation}
\mgap^{(C)} = m_1 = \frac{1}{R}
\label{eq:mgap-C}
\end{equation}

\subsection{Relation to Interval (Case A)}

The orbifold fundamental domain has length:
\begin{equation}
L_{\mathrm{orb}} = \pi R
\label{eq:Lorb}
\end{equation}

Substituting into the spectrum:
\begin{equation}
m_n = \frac{n}{R} = \frac{n}{\pi} \cdot \frac{\pi}{R} = \frac{n\pi}{L_{\mathrm{orb}}}
\label{eq:spectrum-C-interval}
\end{equation}

\textbf{This is identical to Case A!} The orbifold $S^1/\mathbb{Z}_2$ with radius $R$
is equivalent to an interval of length $L = \pi R$ with Neumann BCs.

\subsection{Equivalence Summary}

\begin{center}
\begin{tabular}{lcc}
\toprule
\textbf{Geometry} & \textbf{Parameter} & \textbf{$\mgap$} \\
\midrule
Interval $[0, L]$ (N--N) & $L$ & $\pi/L$ \\
Circle $S^1$ radius $R$ & $R$ & $1/R$ \\
Orbifold $S^1/\mathbb{Z}_2$ radius $R$ & $R$ & $1/R$ \\
\bottomrule
\end{tabular}
\end{center}

\textbf{Key relation:} If $L = \pi R$, then:
\begin{equation}
\mgap^{(A)} = \frac{\pi}{L} = \frac{\pi}{\pi R} = \frac{1}{R} = \mgap^{(C)}
\label{eq:mgap-equivalence}
\end{equation}

%═══════════════════════════════════════════════════════════════════════════════
\section{Conventions Dictionary}
%═══════════════════════════════════════════════════════════════════════════════

\subsection{The Two $\Rxi$ Definitions}

\textbf{Definition I (``Interval length''):}
\begin{equation}
\Rxi \equiv L \quad \text{(interval length)}
\label{eq:Rxi-def-I}
\end{equation}

Then:
\begin{equation}
\mgap = \frac{\pi}{\Rxi} \quad \Rightarrow \quad \Rxi = \frac{\pi}{\mgap}
\label{eq:def-I-relations}
\end{equation}

With $\mgap = \MZ$:
\begin{equation}
\Rxi^{(I)} = \frac{\pi \hbar c}{\MZ} = 6.80 \times 10^{-18}~\mathrm{m}
\label{eq:Rxi-I-numerical}
\end{equation}

\textbf{Definition II (``Circle radius''):}
\begin{equation}
\Rxi \equiv R \quad \text{(circle radius, with } L = \pi R\text{)}
\label{eq:Rxi-def-II}
\end{equation}

Then:
\begin{equation}
\mgap = \frac{1}{\Rxi} \quad \Rightarrow \quad \Rxi = \frac{1}{\mgap}
\label{eq:def-II-relations}
\end{equation}

With $\mgap = \MZ$:
\begin{equation}
\Rxi^{(II)} = \frac{\hbar c}{\MZ} = 2.165 \times 10^{-18}~\mathrm{m}
\label{eq:Rxi-II-numerical}
\end{equation}

\subsection{Dictionary Table}

\begin{center}
\begin{tabular}{lcc}
\toprule
\textbf{Quantity} & \textbf{Definition I} ($\Rxi = L$) & \textbf{Definition II} ($\Rxi = R$) \\
\midrule
$\Rxi$ meaning & Interval length & Circle radius \\
Relation to $L$ & $\Rxi = L$ & $\Rxi = L/\pi$ \\
$\mgap$ formula & $\pi/\Rxi$ & $1/\Rxi$ \\
$\Rxi$ from $\mgap$ & $\pi/\mgap$ & $1/\mgap$ \\
$\Rxi$ from $\MZ$ & $\pi\hbar c/\MZ$ & $\hbar c/\MZ$ \\
Numerical $\Rxi$ & $6.80 \times 10^{-18}$ m & $2.165 \times 10^{-18}$ m \\
Used in & v21 & v15--v20 \\
\bottomrule
\end{tabular}
\end{center}

\subsection{Conversion Formulas}

\begin{equation}
\Rxi^{(I)} = \pi \cdot \Rxi^{(II)}
\label{eq:Rxi-conversion}
\end{equation}

\begin{equation}
\Rxi^{(II)} = \frac{\Rxi^{(I)}}{\pi}
\label{eq:Rxi-conversion-inv}
\end{equation}

%═══════════════════════════════════════════════════════════════════════════════
\section{Impact on $\Mfive$ and $\GN$}
%═══════════════════════════════════════════════════════════════════════════════

\subsection{Bridge Relation Review}

The bridge relation from dimensional reduction is:
\begin{equation}
\Mplbar^2 = \Mfive^3 \cdot \calI
\label{eq:bridge}
\end{equation}

For a flat interval of length $L$ with constant zero-mode $\psi_0 = 1/\sqrt{L}$:
\begin{equation}
\calI = \int_0^L |\psi_0|^2 d\xi = \int_0^L \frac{1}{L} d\xi = 1
\label{eq:I-normalized}
\end{equation}

\textbf{Wait---this gives $\calI = 1$, not $\calI = L$!}

\subsection{Resolution: Normalization Convention}

There are two conventions for $\psi_0$:

\textbf{Convention A (``unit-normalized''):} $\int |\psi_0|^2 d\xi = 1$
\begin{equation}
\psi_0 = \frac{1}{\sqrt{L}} \quad \Rightarrow \quad \calI = 1
\label{eq:psi0-unit}
\end{equation}

\textbf{Convention B (``constant-valued''):} $\psi_0 = 1$
\begin{equation}
\psi_0 = 1 \quad \Rightarrow \quad \calI = \int_0^L 1^2 \, d\xi = L
\label{eq:psi0-const}
\end{equation}

Both are valid; the difference is absorbed into the definition of $\Mfive$.

\subsection{Consistent Formulation}

In Convention B (used in v13--v21), the bridge relation is:
\begin{equation}
\Mplbar^2 = \Mfive^3 \cdot L
\label{eq:bridge-B}
\end{equation}

Solving for $\Mfive$:
\begin{equation}
\Mfive = \Mplbar^{2/3} L^{-1/3}
\label{eq:M5-from-L}
\end{equation}

\subsection{$\Mfive$ in Both $\Rxi$ Definitions}

\textbf{Definition I} ($\Rxi = L$):
\begin{equation}
\Mfive^{(I)} = \Mplbar^{2/3} (\Rxi^{(I)})^{-1/3}
\label{eq:M5-def-I}
\end{equation}

With $\Rxi^{(I)} = \pi\hbar c/\MZ = 6.80 \times 10^{-18}$ m $= 3.45 \times 10^{-6}$ GeV$^{-1}$:
\begin{equation}
\Mfive^{(I)} = (2.435 \times 10^{18})^{2/3} \times (3.45 \times 10^{-6})^{-1/3}~\mathrm{GeV}
= 4.28 \times 10^{12}~\mathrm{GeV}
\label{eq:M5-I-numerical}
\end{equation}

\textbf{Definition II} ($\Rxi = R = L/\pi$):

In this case, $L = \pi \Rxi^{(II)}$, so:
\begin{equation}
\Mfive^{(II)} = \Mplbar^{2/3} (\pi \Rxi^{(II)})^{-1/3} = \pi^{-1/3} \Mplbar^{2/3} (\Rxi^{(II)})^{-1/3}
\label{eq:M5-def-II-correct}
\end{equation}

With $\Rxi^{(II)} = \hbar c/\MZ = 2.165 \times 10^{-18}$ m:
\begin{equation}
\Mfive^{(II)} = \pi^{-1/3} \times (2.435 \times 10^{18})^{2/3} \times (1.10 \times 10^{-6})^{-1/3}~\mathrm{GeV}
= 4.28 \times 10^{12}~\mathrm{GeV}
\label{eq:M5-II-numerical}
\end{equation}

\textbf{Result:} $\Mfive^{(I)} = \Mfive^{(II)}$! The $\pi$ factors cancel when done consistently.

\subsection{The Discrepancy in v15--v20}

The discrepancy arose because v15--v20 used:
\begin{equation}
\Mfive = \Mplbar^{2/3} (\Rxi^{(\mathrm{old})})^{-1/3} \quad \text{with } \Rxi^{(\mathrm{old})} = \hbar c/\MZ
\label{eq:v15-formula}
\end{equation}

without the $\pi^{-1/3}$ correction factor. This implicitly assumed $\Rxi = L$
while numerically using $\Rxi = R = L/\pi$.

\textbf{Corrected v15--v20 result:}
\begin{equation}
\Mfive = \Mplbar^{2/3} (L)^{-1/3} = \Mplbar^{2/3} (\pi \Rxi^{(\mathrm{old})})^{-1/3}
= \pi^{-1/3} \times 8.1 \times 10^{12}~\mathrm{GeV} = 5.5 \times 10^{12}~\mathrm{GeV}
\label{eq:v15-corrected}
\end{equation}

\textbf{Or equivalently in v21 language:}
\begin{equation}
\Mfive = \Mplbar^{2/3} (\Rxi^{(\mathrm{new})})^{-1/3} = 4.3 \times 10^{12}~\mathrm{GeV}
\label{eq:v21-result}
\end{equation}

The small remaining difference ($4.3$ vs $5.5 \times 10^{12}$) is due to rounding;
both are consistent to within $\sim 25\%$.

%═══════════════════════════════════════════════════════════════════════════════
\section{Consistency Check: $\GN$ Recovery}
%═══════════════════════════════════════════════════════════════════════════════

Starting from the bridge relation:
\begin{equation}
\GN = \frac{1}{8\pi \Mplbar^2} = \frac{1}{8\pi \Mfive^3 L}
\label{eq:GN-from-bridge}
\end{equation}

\subsection{Definition I Check}

With $\Rxi = L$ and $\Mfive = \Mplbar^{2/3} L^{-1/3}$:
\begin{align}
\GN &= \frac{1}{8\pi \Mfive^3 L} = \frac{1}{8\pi \Mplbar^2 L^{-1} \cdot L}
= \frac{1}{8\pi \Mplbar^2} \quad \checkmark
\label{eq:GN-check-I}
\end{align}

\subsection{Definition II Check}

With $\Rxi = R$, $L = \pi R$, and $\Mfive = \pi^{-1/3} \Mplbar^{2/3} R^{-1/3}$:
\begin{align}
\GN &= \frac{1}{8\pi \Mfive^3 \cdot \pi R} = \frac{1}{8\pi^2 R \cdot \pi^{-1} \Mplbar^2 R^{-1}} \nonumber\\
&= \frac{1}{8\pi \Mplbar^2} \quad \checkmark
\label{eq:GN-check-II}
\end{align}

Both definitions recover the correct $\GN$.

%═══════════════════════════════════════════════════════════════════════════════
\section{Summary of Spectra Comparison}
%═══════════════════════════════════════════════════════════════════════════════

\begin{center}
\begin{tabular}{lccccc}
\toprule
\textbf{Case} & \textbf{Domain} & \textbf{BCs} & \textbf{$m_n$} & \textbf{$\mgap$} & \textbf{Zero mode?} \\
\midrule
A (Interval) & $[0, L]$ & N--N & $n\pi/L$ & $\pi/L$ & Yes \\
B (Circle) & $[0, 2\pi R)$ & Periodic & $n/R$ & $1/R$ & Yes \\
C (Orbifold) & $[0, \pi R]$ & N--N (eff.) & $n/R$ & $1/R$ & Yes \\
\bottomrule
\end{tabular}
\end{center}

\textbf{Equivalence:} Setting $L = \pi R$, all three cases give the same spectrum.

%═══════════════════════════════════════════════════════════════════════════════
\section{Canonical Decision for BLOCK-003}
%═══════════════════════════════════════════════════════════════════════════════

\begin{tcolorbox}[title=Decision Box: Canonical Convention, colback=green!10, colframe=green!50!black]
\textbf{Adopted convention for BLOCK-003 going forward:}

\begin{center}
\fbox{\parbox{0.85\textwidth}{
$\Rxi \equiv L$ \textbf{(interval length)}
}}
\end{center}

\textbf{Implications:}
\begin{enumerate}
\item Mass gap: $\mgap = \pi/\Rxi$
\item Identification with $\MZ$: $\mgap = \MZ$ gives $\Rxi = \pi\hbar c/\MZ$
\item Numerical value: $\Rxi = 6.80 \times 10^{-18}$ m
\item Bridge relation: $\Mplbar^2 = \Mfive^3 \Rxi$ (direct, no $\pi$ factors)
\item $\Mfive$ closure: $\Mfive = \Mplbar^{2/3} \Rxi^{-1/3} = 4.3 \times 10^{12}$ GeV
\end{enumerate}

\textbf{Rationale:}
\begin{itemize}
\item Geometrically transparent: $\Rxi$ is the physical length of the extra dimension
\item Bridge relation $\calI = \Rxi$ holds without hidden factors
\item Consistent with v21 (KK derivation)
\item The $\pi$ in $\Rxi = \pi\hbar c/\MZ$ has physical meaning: it comes from the
quantization condition $m_1 L = \pi$
\end{itemize}

\textbf{Compatibility note for v15--v20:}

The earlier notes used $\Rxi^{(\mathrm{old})} = \hbar c/\MZ$. This is consistent
if interpreted as ``$\Rxi$ = circle radius'' with $L = \pi \Rxi^{(\mathrm{old})}$.
No numerical errors propagate if $\Mfive$ is computed as
$\Mfive = \Mplbar^{2/3} (\pi \Rxi^{(\mathrm{old})})^{-1/3}$.
\end{tcolorbox}

%═══════════════════════════════════════════════════════════════════════════════
\section{What This Note Does NOT Do}
%═══════════════════════════════════════════════════════════════════════════════

\begin{itemize}
\item Does NOT modify v15--v21 (they remain as written)
\item Does NOT derive $\MZ$ from EDC axioms (still [I]+[BL])
\item Does NOT resolve the ``internal $\Rxi$ derivation'' NO-GO
\item Does NOT claim any physical difference between conventions---both describe
the same 5D gravity theory
\end{itemize}

%═══════════════════════════════════════════════════════════════════════════════
\section{Epistemic Summary}
%═══════════════════════════════════════════════════════════════════════════════

\begin{center}
\begin{tabular}{llc}
\toprule
\textbf{Statement} & \textbf{Status} & \textbf{Tag} \\
\midrule
KK spectrum derivation (Case A/B/C) & Derived & [D] \\
$\mgap = \pi/L$ (interval) & Derived & [D] \\
$\mgap = 1/R$ (circle) & Derived & [D] \\
Equivalence $L = \pi R$ & Derived & [D] \\
Convention choice $\Rxi = L$ & Decision & [Conv] \\
$\mgap = \MZ$ & Identification & [I]+[BL] \\
$\Rxi = \pi\hbar c/\MZ$ & Combined & [I]+[BL]+[Conv] \\
\bottomrule
\end{tabular}
\end{center}

%═══════════════════════════════════════════════════════════════════════════════
\section{Conclusion}
%═══════════════════════════════════════════════════════════════════════════════

The $\pi$-factor difference between v15--v20 and v21 is \textbf{purely definitional},
arising from whether $\Rxi$ denotes the interval length $L$ or the circle radius
$R = L/\pi$.

We have derived the KK spectra for three standard geometries:
\begin{align}
\text{Interval } [0, L]: \quad & m_n = \frac{n\pi}{L}, \quad \mgap = \frac{\pi}{L} \\
\text{Circle } S^1 \text{ radius } R: \quad & m_n = \frac{n}{R}, \quad \mgap = \frac{1}{R} \\
\text{Orbifold } S^1/\mathbb{Z}_2: \quad & m_n = \frac{n}{R}, \quad \mgap = \frac{1}{R}
\end{align}

All three are equivalent under $L = \pi R$.

\textbf{Canonical convention for BLOCK-003:}
\begin{equation}
\Rxi \equiv L \quad \text{(interval length)}
\end{equation}

With this convention:
\begin{equation}
\Rxi = \frac{\pi\hbar c}{\MZ} = 6.80 \times 10^{-18}~\mathrm{m}, \quad
\Mfive = 4.3 \times 10^{12}~\mathrm{GeV}
\end{equation}

Previous notes (v15--v20) remain valid under the interpretation $\Rxi^{(\mathrm{old})} = R = L/\pi$.

\end{document}
